\newpage
\section{\tkzClass{quadrilateral} 类}

\vspace{1em}
变量 \tkzVar{quadrilateral}{Q} 持有一个表，用于存放四边形。它是可选的，你可以自由选择变量名。不过，为清晰和一致起见，推荐使用 \code{Q}。如果你使用自定义变量（例如 Quad），则必须手动初始化它。如果使用了 \code{Q}，\tkzFct{tkz-elements}{init\_elements()} 函数会重新初始化该表。


\subsection{创建四边形}
\label{sub:creating_a_quadrilateral}
\tkzClass{quadrilateral} 类需要四个点。点的顺序决定四边形的边。

\medskip
该对象通常存放在 \tkzVar{quadrilateral}{Q} 中，这是推荐的变量名。

\begin{mybox}
\begin{verbatim}
Q.ABCD = quadrilateral:new(z.A, z.B, z.C, z.D)
\end{verbatim}
\end{mybox}

\textbf{简写形式。}

也提供了简写构造器：

\begin{mybox}
\begin{verbatim}
Q.ABCD = quadrilateral(z.A, z.B, z.C, z.D)
\end{verbatim}
\end{mybox}

\subsection{四边形的属性}
四个点按顺序给出；不做检验。

\begin{center}
 \code{ 创建：Q.new = quadrilateral(z.A,z.B,z.C,z.D)}
\end{center}


  \bgroup
  \catcode`_=12
  \small
  \captionof{table}{四边形的属性。}\label{quadrilateral:attributes}
  \begin{tabular}{lll}
  \toprule
  \textbf{属性}         & \textbf{用法} & \\
  \tkzAttr{quadrilateral}{pa}   & |z.A = Q.new.pa| & \\
  \tkzAttr{quadrilateral}{pb}   & |z.B = Q.new.pb| & \\
  \tkzAttr{quadrilateral}{pc}   & |z.C = Q.new.pc| & \\
  \tkzAttr{quadrilateral}{pd}   & |z.D = Q.new.pd| & \\
  \tkzAttr{quadrilateral}{type} & |Q.new.type= 'quadrilateral'|  &\\
  \tkzAttr{quadrilateral}{center}    & |z.I = Q.new.center| & 对角线交点\\
  \tkzAttr{quadrilateral}{g}    & |z.G = Q.new.g| & 重心\\
  \tkzAttr{quadrilateral}{a}    & |AB = Q.new.a| & AB 的长度\\
  \tkzAttr{quadrilateral}{b}    & |BC = Q.new.b| & BC 的长度\\
  \tkzAttr{quadrilateral}{c}    & |CD = Q.new.c| & CD 的长度\\
  \tkzAttr{quadrilateral}{d}    & |DA = Q.new.d| & DA 的长度\\
  \tkzAttr{quadrilateral}{ab}   & |Q.new.ab|   &  过两个顶点的直线   \\
  \tkzAttr{quadrilateral}{ac}   & |Q.new.ac|   &  同上。 \\
  \tkzAttr{quadrilateral}{da}   & |Q.new.da|   &  同上。 \\
  \tkzAttr{quadrilateral}{bc}   & |Q.new.bc|   &  同上。 \\
  \tkzAttr{quadrilateral}{bd}   & |Q.new.bd|   &  同上。 \\
  \tkzAttr{quadrilateral}{cd}   & |Q.new.cd|   &  同上。 \\
  \bottomrule
  \end{tabular}
  \egroup


\subsubsection{四边形的属性}

\vspace{1em}
\begin{minipage}{.5\textwidth}
\begin{verbatim}
\directlua{
init_elements()
z.A = point(0, 0)
z.B = point(4, 0)
z.C = point(5, 1)
z.D = point(0, 3)
Q.ABCD = quadrilateral(z.A, z.B, z.C, z.D)
z.I = Q.ABCD.i
z.G = Q.ABCD.g}

\begin{tikzpicture}
\tkzGetNodes
\tkzDrawPolygon(A,B,C,D)
\tkzDrawSegments(A,C B,D)
\tkzDrawPoints(A,B,C,D,I,G)
\end{tikzpicture}
\end{verbatim}
\end{minipage}
\begin{minipage}{.5\textwidth}
\directlua{
 init_elements()
 z.A = point(0, 0)
 z.B = point(4, 0)
 z.C = point(5, 1)
 z.D = point(0, 3)
 Q.ABCD = quadrilateral(z.A, z.B, z.C, z.D)
 z.I = Q.ABCD.center
 z.G = Q.ABCD.g}

\begin{center}
  \begin{tikzpicture}
  \tkzGetNodes
  \tkzDrawPolygon(A,B,C,D)
  \tkzDrawSegments(A,C B,D)
  \tkzDrawPoints(A,B,C,D,I,G)
  \end{tikzpicture}
\end{center}
\end{minipage}

\subsubsection{四边形示例}
高等欧几里得几何 2013 年 6 月 26 日补充材料\\

\texttt{几何平均的构造：}


\directlua{
init_elements()
z.A = point(0, 0)
z.B = point(1, 4)
z.C = point(3.5, 4)
z.D = point(5, 0 )
Q.ABCD = quadrilateral(z.A, z.B, z.C, z.D)
z.I = Q.ABCD.center
L.ll = Q.ABCD.da:ll_from(z.I)
z.H = intersection(L.ll, Q.ABCD.ab)
z.Hp = intersection(L.ll, Q.ABCD.cd)
z.M =  Q.ABCD.ab.mid
C.MA = circle(z.M, z.A)
L.perp = Q.ABCD.ab:ortho_from(z.H)
_,z.P = intersection (C.MA, L.perp)
T.PAB = triangle (z.P, z.A, z.B)
L.bis =  T.PAB:bisector ()
z.G = L.bis.pb
L.llg = Q.ABCD.da:ll_from(z.G)
L.llm = Q.ABCD.da:ll_from(z.M)
z.Gp = intersection(L.llg, Q.ABCD.cd)}
\begin{center}
\begin{tikzpicture}
\tkzGetNodes
\tkzDrawPolygon(A,B,C,D)
\tkzDrawCircle(M,A)
\tkzDrawSegments[red](H,H')
\tkzDrawSegments[blue](G,G')
\tkzDrawSegments[dashed](P,H P,B P,A P,G A,C B,D)
\tkzDrawPoints(A,B,C,D,H,I,M,P,G,G',H')
\tkzLabelPoints(A,D,I)
\tkzLabelPoints[above](B,C,P)
\tkzLabelPoints[above right=-3pt](H,G)
\tkzLabelPoints[right](G',H')
 \tkzLabelSegment[above](B,C){$b$}
 \tkzLabelSegment[below](A,D){$a$}
\end{tikzpicture}
\end{center}



\begin{proof}
给定一个梯形，其平行边为 $AD$ 和 $BC$。

设 H 为过 $I$ 且平行于边 $AD$ 的直线的端点。则 $HH' = \dfrac{2ab}{a+b}$。

$PH$ 垂直于 $AB$

$(PG)$ 是 $\widehat{APB}$ 的角平分线

 $GG' = \dfrac{a \cdot \sqrt{b} + b \cdot \sqrt{a}}{\sqrt{a}+\sqrt{b}}
= \dfrac{\sqrt{ab}(\sqrt{a}+\sqrt{b})}{\sqrt{a}+\sqrt{b}} = \sqrt{ab} $
\end{proof}

\texttt{代码：}

\begin{verbatim}
\directlua{
init_elements()
z.A = point(0, 0)
z.B = point(1, 4)
z.C = point(3.5, 4)
z.D = point(5, 0 )
Q.ABCD = quadrilateral(z.A, z.B, z.C, z.D)
z.I = Q.ABCD.center
L.ll = Q.ABCD.da:ll_from (z.I)
z.H = intersection(L.ll, Q.ABCD.ab)
z.Hp = intersection(L.ll, Q.ABCD.cd)
z.M =  Q.ABCD.ab.mid
C.MA = circle(z.M,z.A)
L.perp = Q.ABCD.ab:ortho_from(z.H)
_,z.P = intersection(C.MA, L.perp)
T.PAB = triangle(z.P, z.A, z.B)
L.bis =  T.PAB:bisector()
z.G = L.bis.pb
L.llg = Q.ABCD.da:ll_from(z.G)
L.llm = Q.ABCD.da:ll_from(z.M)
z.Gp = intersection(L.llg, Q.ABCD.cd)}
\end{verbatim}

\subsection{四边形的方法}

\vspace{1em}

  \bgroup
  \catcode`_=12
  \small
  \captionof{table}{四边形的方法。}\label{quadrilateral:methods}
  \begin{tabular}{ll}
  \toprule
  \textbf{方法} & \texttt{参考}    \\
  \midrule
  \tkzMeth{quadrilateral}{new()}& [\ref{sub:creating_a_quadrilateral}]\\

  \tkzMeth{quadrilateral}{is\_cyclic ()} & [\ref{ssub:iscyclic}]\\
  \tkzMeth{quadrilateral}{is\_convex ()} & [\ref{ssub:isconvex}]\\
  \tkzMeth{quadrilateral}{poncelet\_point()} & [\ref{ssub:poncelet_point}]\\
  \tkzMeth{quadrilateral}{area()} & [\ref{ssub:quadrilateral_area}]\\
  \tkzMeth{quadrilateral}{get(<i>)} & [\ref{ssub:method_quadrilateral_get}]\\
  \bottomrule %
  \end{tabular}
  \egroup



\subsubsection{方法 \tkzMeth{quadrilateral}{get(<i>)}}
\label{ssub:method_quadrilateral_get}

该方法返回定义四边形对象的各个点。

\begin{itemize}
  \item \code{Q.ABCD:get()} 按顺序返回四个顶点：\code{pa}、\code{pb}、\code{pc}、\code{pd}；
  \item \code{Q.ABCD:get(i)}（\code{i} 取 1 到 4）只返回第 \code{i} 个顶点；例如 \code{Q.ABCD:get(2)} 等价于 \code{Q.ABCD.pb}。
\end{itemize}

\begin{tkzexample}[code only]
\directlua{
  init_elements()
  z.A = point(0, 0)
  z.B = point(2, 0)
  z.C = point(3, 2)
  z.D = point(0, 1)
  Q.ABCD = quadrilateral(z.A, z.B, z.C, z.D)
  z.E, z.F, z.G, z.H = Q.ABCD:get()
  z.K = Q.ABCD:get(3)
}
\end{tkzexample}

\subsubsection{方法 \tkzMeth{quadrilateral}{is\_cyclic()}}
\label{ssub:iscyclic}

若四个顶点共圆（容差为 \code{tkz.epsilon}）则返回 \code{true}，否则返回 \code{false}。

\textbf{圆内接四边形}

\vspace{1em}
\begin{minipage}{.5\textwidth}
\begin{verbatim}
\directlua{
 init_elements()
 z.A = point(0, 0)
 z.B = point(4, 0)
 z.D = point:polar(4, tkz.tau / 3)
 L.DB = line(z.D, z.B)
 T.equ = L.DB:equilateral()
 z.C = T.equ.pc
 Q.new = quadrilateral(z.A, z.B, z.C, z.D)
 bool = Q.new:is_cyclic()
 if bool == true then
 	C.cir = triangle(z.A, z.B, z.C):
 	 circum_circle()
 	z.O = C.cir.center
 end}
\begin{tikzpicture}
  \tkzGetNodes
  \tkzDrawPolygon(A,B,C,D)
  \tkzDrawPoints(A,B,C,D)
  \tkzDrawCircle(O,A)
  \ifthenelse{\equal{\tkzUseLua{bool}}{
     true}}{\tkzDrawCircle(O,A)}{}
  \tkzLabelPoints(A,B)
  \tkzLabelPoints[above](C,D)
\end{tikzpicture}
\end{verbatim}
\end{minipage}
\begin{minipage}{.5\textwidth}
\directlua{
init_elements()
z.A = point(0, 0)
z.B = point(4, 0)
z.D = point:polar(4, tkz.tau / 3)
L.DB = line(z.D, z.B)
T.equ = L.DB:equilateral()
z.C = T.equ.pc
Q.new = quadrilateral(z.A, z.B, z.C, z.D)
bool = Q.new:is_cyclic()
if bool == true then
	C.cir = triangle(z.A, z.B, z.C):circum_circle()
	z.O = C.cir.center
end
}
\begin{center}
\begin{tikzpicture}[scale = .75]
  \tkzGetNodes
  \tkzDrawPolygon(A,B,C,D)
  \tkzDrawPoints(A,B,C,D)
  \tkzDrawCircle(O,A)
  \ifthenelse{\equal{\tkzUseLua{bool}}{
   true}}{\tkzDrawCircle(O,A)}{}
  \tkzLabelPoints(A,B)
  \tkzLabelPoints[above](C,D)
\end{tikzpicture}
\end{center}
\end{minipage}

\subsubsection{方法 \tkzMeth{quadrilateral}{is\_convex()}}
\label{ssub:isconvex}

若四边形为凸则返回 \code{true}，否则返回 \code{false}。三个连续顶点共线的四边形不是凸的。

\begin{verbatim}
if Q.ABCD:is_convex() then ... end
\end{verbatim}

\subsubsection{方法 \tkzMeth{quadrilateral}{area()}}
\label{ssub:quadrilateral_area}

返回四边形的面积，无论是否为凸。对于自交四边形会引发错误。

\begin{verbatim}
local s = Q.ABCD:area()
\end{verbatim}







\subsubsection{方法 \tkzMeth{quadrilateral}{poncelet\_point}}
\label{ssub:poncelet_point}

定义见 [\ref{ssub:method_triangle_poncelet_point}]。
Poncelet 点对任意四边形都有定义；若有三个顶点共线则引发错误。

\begin{tkzexample}[latex=.5\textwidth]
\directlua{
 init_elements()
 z.A = point(1, 1)
 z.B = point(6, 0)
 z.D = point(0, 5)
 z.C = point(6, 2)
 Q.ABCD = quadrilateral(z.A, z.B, z.C, z.D)
 z.P = Q.ABCD:poncelet_point()
 T.ABC = triangle(z.A,z.B,z.C)
 z.I = T.ABC.eulercenter
 z.Mc = tkz.midpoint(z.A, z.B)
 T.ABD = triangle(z.A, z.B, z.D)
 z.I1 = T.ABD.eulercenter}
\begin{center}
\begin{tikzpicture}
  \tkzGetNodes
  \tkzDrawPolygons(A,B,C,D)
  \tkzDrawPoints(A,B,C,D,P,Mc)
  \tkzDrawCircles[red](I,Mc I1,Mc)
  \tkzLabelPoints(A,B)
  \tkzLabelPoints[above](C,D,P)
\end{tikzpicture}
\end{center}
\end{tkzexample}
\endinput